%% IAU: Assisting the CPU Core for Next-Generation Multi-Million IOPS NVMe Devices
%% Target venue: IEEE Computer Architecture Letters (CAL)
%% Format: IEEEtran, two-column, 10pt; 4 pages text + ~1 page refs
%%
%% Drafting status (2026-05-11, third revision):
%%   - Abstract  : Tight 4-move structure (background / problem / design /
%%                 effectiveness). Numbers left as TODO placeholders until
%%                 Mode 2A/2B sweeps land.
%%   - Section 1 : fully drafted; references [1]-[5], [8], [18] wired in.
%%   - Section 2 : Mode 0 QD-sweep numbers (QD=16/32/64/128) wired in for the
%%                 per-IO decomposition and the polling-asymptote anchor stat.
%%                 State_Dealloc split still pending TODO[split_rebake].
%%   - Sections 3-6: not yet drafted in this pass.
%%
%% TODO markers still in use:
%%   TODO[baseline]            -- Mode 0 baseline IOPS at QD=128 in Abstract
%%   TODO[iau]                 -- IAU-mode IOPS at QD=128 in Abstract
%%   TODO[speedup]             -- IAU / baseline throughput ratio in Abstract
%%   TODO[cycles_pct]          -- per-IO cycle reduction percent in Abstract
%%   TODO[mode_bracket]        -- Mode 2A / Mode 2B savings (elsewhere in body)
%%   TODO[split_rebake]        -- State_Dealloc library/callback split
%%   TODO[author_info]         -- fill in author/affiliation block
%%
\documentclass[journal]{IEEEtran}

\usepackage{cite}
\usepackage{amsmath,amssymb,amsfonts}
\usepackage{graphicx}
\usepackage{textcomp}
\usepackage{xcolor}
\usepackage{booktabs}
\usepackage{url}

% Compact spacing helpers; CAL has a hard 4-page limit on the body.
\usepackage[compact]{titlesec}

\begin{document}

% ============================================================================
% Title block
% ============================================================================
\title{IAU: Assisting the CPU Core for Next-Generation\\Multi-Million IOPS NVMe Devices}

\author{%
  % TODO[author_info]: replace placeholders below.
  Anonymous Authors%
  \thanks{Manuscript submitted to IEEE Computer Architecture Letters,
  \today.}%
}

\markboth{IEEE Computer Architecture Letters, Vol.~XX, No.~Y, 20XX}%
{Author~\MakeLowercase{\textit{et~al.}}: IAU}

\maketitle

% ============================================================================
% Abstract  --  Mode 0 baseline wired (QD=128, 4 KB random read).
%                TODO[mode_bracket]: insert Mode 2A / Mode 2B savings.
% ============================================================================
\begin{abstract}
Next-generation NVMe SSDs are reaching tens of millions of IOPS to serve
retrieval-augmented inference, billion-scale vector search, and
graph workloads that treat NAND as an active memory tier rather than a
bulk store. At that rate, the host CPU---not the device---becomes the
bottleneck: even a state-of-the-art user-space polling stack spends the
bulk of its cycles on per-IO bookkeeping, address translation, and
empty completion-queue scans, leaving the device chronically
under-utilized. We propose \emph{IAU} (I/O Assistant Uncore), an
on-chiplet block in the CPU uncore---placed alongside the PCIe root
complex and integrated memory controller---that executes NVMe queues in
hardware via SRAM-resident submission/completion queues, doorbell
aggregation, and completion suppression, while preserving the existing
SPDK programming model. On a full-system gem5+SimpleSSD+SPDK platform,
IAU lifts sustained 4\,KB random-read throughput from
TODO[baseline]~M~IOPS to TODO[iau]~M~IOPS at QD=128---a
TODO[speedup]$\times$ improvement---while reducing per-IO host cycles
by TODO[cycles\_pct]\,\% within a per-block SRAM footprint of 329~KB to
2.6~MB synthesized for ASAP7 7\,nm.
\end{abstract}

\begin{IEEEkeywords}
NVMe, SSD, storage stacks, CPU uncore, chiplet architecture, queue
offload, million-IOPS storage, full-system simulation, SPDK.
\end{IEEEkeywords}

% ============================================================================
% Section 1 -- Introduction
%
% Locked structure (PAPER_CHAPTER_PLAN.md):
%   Para 1: HOOK + WORKLOAD          refs [1][2][5]
%   Para 2: WALL                     refs [3][18]
%   Para 3: PLACEMENT + ANALOGY      refs [4][8]
%   Para 4: CONTRIBUTIONS            (no refs)
% Both memorable lines wired in verbatim.
% ============================================================================
\section{Introduction}
\label{sec:intro}

\IEEEPARstart{T}{he} arrival of solid-state devices capable of sustaining
$10^8$ random IOPS, demonstrated publicly in 2025 alongside roadmaps from
storage vendors targeting GPU-storage direct paths for inference
serving~\cite{kioxia_fms2025}, has shifted what storage-touching workloads
look like. In-storage retrieval-augmented generation pipelines now spend the
majority of their wall-clock time on storage access: a recent in-storage
acceleration study reports that 61\,\% of the cost of retrieval-augmented
generation as a service is the retrieval read path itself~\cite{ragx_isca25}.
Underneath these workloads, GPU-initiated I/O frameworks such as BaM submit
NVMe requests directly from streaming multiprocessors and consume billions
of small reads per second on a single
node~\cite{bam_asplos23}. The data device side has thus moved from
\emph{passive bulk store} to \emph{active memory tier}: NAND-resident
indices and embeddings are read with the access pattern and rate of DRAM,
not files. AI retrieval workloads are turning NAND into an active memory
tier, but the host CPU is the wall.

That wall is software. State-of-the-art user-space NVMe drivers replace
interrupts with polling, replace per-IO system calls with shared
submission/completion queues, and re-use pinned huge-page buffers to keep
the IO path off the kernel boundary~\cite{spdk_modern_apis,spdk_docs}. Even
so, every IO incurs the same hot loop on the host core: tracker allocation,
PRP-list page-walks, command construction, a fenced doorbell ring, a
DRAM-resident completion-queue scan, and a callback invocation. The result
is a per-IO software floor that the recent measurement-based survey of
\emph{libaio}, \emph{io\_uring}, and SPDK independently
confirmed~\cite{spdk_modern_apis}. When the device can supply tens of
millions of IOPS, that floor becomes the binding constraint: on our
gem5+SimpleSSD+SPDK platform with a 32~M~IOPS device ceiling, unmodified
SPDK at QD=128 delivers 0.82~M~IOPS (a $\sim$40\,$\times$ shortfall) while
\emph{idling} on 97\,\% of its polling calls---no amount of batching,
polling tuning, or NUMA pinning crosses it, because the floor is set by
memory-hierarchy round-trips and an $\mathcal{O}$(qpairs) completion scan
that neither queue-depth nor core-count can amortize.

The historical resolution to a similar architectural wall is informative.
When DRAM access latencies started to dominate processor pipelines in the
early 2000s, AMD relocated the memory controller from the chipset onto the
CPU die, eliminating a fabric hop and changing the latency
class~\cite{opteron_imc}. That decision was not about adding a faster
external memory; it was about placement. The cost-benefit framing of which
storage tier deserves what level of integration was independently re-derived
for modern hardware by the thirty-year retrospective on the
five-minute-rule~\cite{five_minute_rule_30}, which places NVMe firmly in
the on-die-integration regime now that device latency is at the
microsecond mark. We close the wall where the memory-controller wall was
closed: in the CPU uncore.

We make the following contributions.

\begin{enumerate}
  \item \textbf{An IAU architecture} that relocates NVMe queue
        execution into the CPU uncore, organized around three
        cooperating mechanisms---SRAM-resident submission/completion
        queues, doorbell aggregation, and completion suppression---and a
        dual software contract: a transparent Mode~2A requiring no SPDK
        changes, and a poll-lite Mode~2B requiring a $\sim$45-line patch
        confined to a single hot-path polling function.
  \item \textbf{A measurement-based characterization} of the per-IO
        software overhead in SPDK on a full-system
        gem5+SimpleSSD+SPDK platform with a 32~M~IOPS device ceiling,
        under which unmodified SPDK saturates at 0.82~M~IOPS at QD=128,
        decomposed into named per-stage costs that expose the
        memory-hierarchy floor and the $\mathcal{O}$(qpairs) completion
        scan as the structural asymptotes.
  \item \textbf{A bracket evaluation} showing that the transparent mode
        recovers the elidable host work under unmodified SPDK (lower
        bound), while the poll-lite mode further eliminates the
        multi-queue scanning overhead with a minimal patch (upper
        bound), with SRAM and area envelopes anchored on ASAP7 7\,nm
        Yosys synthesis for four design points spanning 329~KB to
        2.6~MB per uncore block.
\end{enumerate}

% ============================================================================
% Section 2 -- The Host I/O Wall
%
% Locked structure (PAPER_CHAPTER_PLAN.md):
%   2.1 Per-IO Cost Decomposition (~250 words + FIG 2)
%   2.2 Why Software Cannot Close This Gap (~400 words)
%       memory-hierarchy floor + multi-queue scanning O(qpairs)
%   2.3 Bridge sentence (~50 words)
%
% Anchor stat (2026-05-11 Mode 0 sweep, 4 KB random read, qpair=1):
%   QD=16   : 776,253 IOPS,  86.0% idle polls, Completions_Per_Call = 2.29
%   QD=32   : 799,927 IOPS,  90.9% idle polls, Completions_Per_Call = 2.91
%   QD=64   : 812,491 IOPS,  94.7% idle polls, Completions_Per_Call = 3.38
%   QD=128  : 819,253 IOPS,  97.1% idle polls, Completions_Per_Call = 3.66
%   IOPS asymptote ~0.82 M; SSD ceiling ~32 M -> ~40x shortfall.
%
% Energy out of scope per ¶§2.2_scope_locked.
% ============================================================================
\section{The Host I/O Wall}
\label{sec:wall}

\subsection{Per-IO cost decomposition}
\label{sec:wall:decomp}

To attribute the gap between device capability and observed throughput, we
instrument SPDK's per-IO path along the eight stages it traverses on every
submission and completion---submit-preamble, tracker allocation, PRP-list
construction (address translation), command construction, fence, doorbell
write, completion-queue-entry phase-bit detection, tracker lookup, and
state deallocation---and report the per-stage cost in nanoseconds against
the cycles-per-IO budget that the device implies. The instrumentation
splits the legacy \emph{state-deallocation} stage into a library-teardown
component (CID free-list splice, tracker push, qpair statistics) and an
application-callback component (workload-defined; for \texttt{spdk\_nvme\_perf}
this is a latency-histogram update). The split matters because hardware
can absorb the library teardown but cannot replace application code;
reporting the two as one stage would credit hardware for cycles it cannot
fundamentally reach.

The instrumented path reveals that the dominant per-IO work is not the
doorbell or fence---both consume single-digit nanoseconds in the
steady-state polling regime (1.06\,ns and 1.82\,ns per IO at QD=128,
respectively)---but a combination of address translation, state
deallocation, and an unattributed per-call overhead that together absorb
roughly two-thirds of the elidable per-IO budget.
Across a 4\,KB random-read Mode~0 sweep on our gem5+SimpleSSD+SPDK
platform, host throughput saturates at \textbf{0.82~M~IOPS} at
QD=128---only 2.6\,\% of the 32~M~IOPS device ceiling---while
\textbf{97.1\,\%} of polling calls return zero useful completions
(i.e., the host scans the completion queue and finds nothing) and the
average useful-completions-per-call is \textbf{3.66}. The asymptote is
visible across the QD sweep: throughput climbs only modestly from
776,253~IOPS at QD=16 to 819,253~IOPS at QD=128 ($+5.5$\,\%), while
median latency grows almost linearly with QD (20.5, 39.9, 78.8,
156.7\,$\mu$s), confirming that work is queueing in the host---not the
device---across the regime.
As a corroborating low-QD check, the QD=16 datum represents a
$4.6\times$ lift over the SSD-cycle-bound baseline of the same
simulator, with an average per-IO software cost of 8.04\,$\mu$s on the
submission side plus 1.42\,$\mu$s on the completion side.
The behavior is monotonic across queue depths: as the device gets out of
the way, the host wastes a larger fraction of cycles on bookkeeping and on
scanning queues for work that has not yet arrived.

Figure~\ref{fig:io-breakdown} (placeholder) decomposes per-IO cycles into
the named stages plus an \emph{other-SPDK-overhead} residual at two queue
depths, distinguishing hardware-elidable cycles from the
application-callback component that any architectural offload must
preserve. At QD=128 the per-IO costs of the largest elidable stages are
140.6\,ns (submit preamble), 217.8\,ns (tracker allocation), 337.0\,ns
(PRP-list address translation), 51\,ns (command construction),
20\,ns (CQE phase-bit detection), 11\,ns (tracker lookup), and 295.2\,ns
(legacy state-deallocation bundle)---roughly 1{,}073\,ns of per-IO
software-fabric work, against $\sim$2.9\,ns of physical-boundary work
(doorbell + fence combined) and a $\sim$1.22\,$\mu$s per-IO budget set
by the 0.82~M~IOPS asymptote. The full per-stage table after the
State\_Dealloc split is populated by TODO[split\_rebake]; the
application-callback portion of the legacy bundle is expected to be on
the order of tens of nanoseconds and the library teardown the remainder.

\subsection{Why software cannot close this gap}
\label{sec:wall:floors}

Two independent floors bound any purely software stack at the IOPS
regime motivated in Section~\ref{sec:intro}.

\paragraph*{Memory-hierarchy floor}
Every IO performs at least three DRAM round-trips that no batching
strategy can eliminate: one read of the submission-queue tracker pool to
allocate a command identifier, one DRAM-resident write of the NVMe
submission queue entry, and one read of the completion queue entry on
the polling side. At a conservative 80\,ns per round-trip and a single
core, these three traversals alone establish a $\sim$240\,ns/IO floor
that bounds aggregate throughput at roughly 4\,M~IOPS per core before
any productive work is performed. Our measured Mode~0 throughput of
0.82~M~IOPS sits a factor of $\sim$5$\times$ below this DRAM-limited
ceiling, which tells us two things: (i) the per-stage software work
characterized in Section~\ref{sec:wall:decomp}---and not memory
bandwidth---is what binds the host today, and (ii) even an oracle
elimination of every elidable software stage cannot push a single core
past the 4\,M~IOPS structural ceiling. Figure~\ref{fig:dram-floor}
(placeholder) reports DRAM bytes per IO from the gem5 memory-system
counters; the measurement is expected to be invariant under queue depth
and core count for fixed IO size, as required of a structural floor.

\paragraph*{Multi-queue scanning $\mathcal{O}(\textrm{qpairs})$}
Modern NVMe deployments use one submission/completion queue pair per
core or per worker thread, and the SPDK polling routine scans each
qpair's completion queue head pointer on every polling call. The cost
of producing a single useful completion therefore scales with the
number of registered qpairs, not with the number of completions
actually waiting. Even at the single-qpair configuration measured here,
the symptom is already visible: 97.1\,\% of polling calls at QD=128
return zero completions, and the
\texttt{scans\_per\_completion} statistic of 0.273 reflects an
amortization advantage that exists only because doorbell aggregation
batches eight to sixteen IOs per ring---an advantage that erodes as
soon as a second qpair must also be scanned on every poll. A multi-qpair
sweep is in progress and will report the linear growth in
\texttt{scans\_per\_completion} and the corresponding fall in
completions-per-call as qpair count rises at fixed device load.

These two floors are structural, not implementation-specific: the
memory-hierarchy floor is set by physics; the multi-queue scan is set
by the NVMe submission model in which the host must visit every queue
to find new entries. Optimizing the per-stage costs that
Section~\ref{sec:wall:decomp} enumerates flattens the curve, but cannot
move the asymptote. The case for hardware is the union of both floors.
\emph{Energy per IO is intentionally out of scope in this section:
per-IO joules track per-IO cycles modulo a constant-energy-per-instruction
factor; reporting it adds no independent information beyond what the
cycles-per-IO measurement above already establishes, and on-die power
feasibility for the proposed mechanism is covered in
Section~\ref{sec:eval} from RTL synthesis.}

\subsection{Bridge}
\label{sec:wall:bridge}

The work above the device is hot-loop bookkeeping that the host CPU is
structurally bad at and which the device cannot help with. The natural
fix is to give it that hardware.

% ============================================================================
% Sections 3-6: NOT YET DRAFTED in this pass.
%
% Placeholders so the document compiles and the section structure is visible.
% Drafting order from PAPER_CHAPTER_PLAN.md section 4 reading guide:
%   sec:arch (3)  -> next pass; depends on Figure 1 and RTL_design reports
%   sec:eval (4)  -> next pass; depends on Mode 1/2 sweep completion
%   sec:related (5) -> next pass; Table 1 already designed
%   sec:conclusion (6) -> last pass; closing line locked
% ============================================================================
\section{IAU Architecture}
\label{sec:arch}
% TODO[sec3]: full draft in the next pass.

\section{Evaluation}
\label{sec:eval}
% TODO[sec4]: full draft once Mode 1 + Mode 2B sweeps complete.

\section{Related Work}
\label{sec:related}
% TODO[sec5]: full draft; Table 1 already designed in PAPER_CHAPTER_PLAN.md.

\section{Conclusion}
\label{sec:conclusion}
% TODO[sec6]: closing line locked: "The uncore is the placement that
% scales."

% ============================================================================
% Bibliography
% ============================================================================
\bibliographystyle{IEEEtran}
\bibliography{paper}

\end{document}
